\section{Phase 9: A Minimal Differentiable Optimization Loop for Raw Scheme C}
\label{sec:phase9_minimal_differentiable_loop}

The purpose of this phase is not yet to construct a full planning framework, but to verify that BLOM already functions as a \emph{differentiable trajectory representation} that can be inserted into an optimization loop.
At the present stage, the only admissible optimization object is the raw Scheme~C map in the canonical setting:
\[
s=4,\qquad k=2.
\]
The optimization loop considered in this phase is therefore deliberately minimal:
\[
\theta=(\bar q,T)
\ \longmapsto\
c(\theta)=c^{\mathrm{C},2}(\bar q,T)
\ \longmapsto\
K(c,T)
\ \longmapsto\
\nabla_\theta J(\theta),
\qquad
J(\theta):=K(c(\theta),T).
\]
The goal is to prove that this chain is mathematically well defined, differentiable, exactly implementable in block-banded form, and capable of producing a descent step under standard smoothness assumptions.

Throughout this section, all notation is inherited from Phases~0--8.
In particular,
\[
q_0,q_1,\dots,q_M\in\mathbb R,
\qquad
\bar q:=(q_1,\dots,q_{M-1})^\top\in\mathbb R^{M-1},
\qquad
T=(T_1,\dots,T_M)^\top\in\mathbb R_{>0}^M,
\]
and
\[
\theta:=(\bar q,T)\in\mathbb R^{M-1}\times \mathbb R_{>0}^{M}.
\]
The raw Scheme~C coefficient vector is
\[
c(\theta)=c^{\mathrm{C},2}(\bar q,T)\in\mathbb R^{2sM}=\mathbb R^{8M},
\]
with segment blocks
\[
c(\theta)=
\begin{bmatrix}
c_1(\theta)\\
\vdots\\
c_M(\theta)
\end{bmatrix},
\qquad
c_i(\theta)\in\mathbb R^{2s}=\mathbb R^8.
\]

\subsection{Optimization scope and admissible parameter domain}
\label{subsec:phase9_scope}

This phase only considers the raw Scheme~C map and excludes:
\begin{itemize}
    \item Scheme~A and Scheme~B assembly,
    \item corrected variants such as C1/C2,
    \item and any full \(k\)-family planner.
\end{itemize}
The admissible parameter domain is
\begin{equation}
\label{eq:phase9_theta_adm}
\Theta_{\mathrm{adm}}(T_{\min})
:=
\Bigl\{
(\bar q,T)\in \mathbb R^{M-1}\times \mathbb R^M
\ \Big|\
T_i>T_{\min}>0,\ \forall i
\Bigr\}.
\end{equation}
The strict positivity margin \(T_{\min}\) is introduced so that local time perturbations remain admissible during optimization.

\subsection{Minimal objective functional}
\label{subsec:phase9_objective}

We define a minimal objective
\begin{equation}
\label{eq:phase9_total_objective}
J(\theta):=
K(c(\theta),T)
=
K_{\mathrm{ctrl}}(c,T)
+
\lambda_T K_{\mathrm{time}}(T)
+
\lambda_{\mathrm{obs}} K_{\mathrm{soft\mbox{-}obs}}(c,T),
\end{equation}
where \(\lambda_T\ge 0\) and \(\lambda_{\mathrm{obs}}\ge 0\) are fixed weights.

\subsubsection{Control cost}
\label{subsubsec:phase9_control_cost}

For each segment \(i\), let \(H_i(T_i)\in\mathbb R^{2s\times 2s}\) be the positive semidefinite segment Gram matrix associated with the minimum-\(s\)-th-derivative cost, as in Phase~1.
Define
\begin{equation}
\label{eq:phase9_Kctrl}
K_{\mathrm{ctrl}}(c,T)
:=
\frac12\sum_{i=1}^{M} c_i^\top H_i(T_i)c_i.
\end{equation}
In the canonical minimum-snap case \(s=4\), \(H_i(T_i)\) is the segment snap Gram matrix.

\subsubsection{Time regularization}
\label{subsubsec:phase9_time_cost}

Define
\begin{equation}
\label{eq:phase9_Ktime}
K_{\mathrm{time}}(T):=\sum_{i=1}^{M}T_i.
\end{equation}

\subsubsection{Soft obstacle penalty}
\label{subsubsec:phase9_soft_obs}

Fix a finite set of normalized sample locations
\[
0\le \alpha_1,\dots,\alpha_{N_s}\le 1,
\qquad
w_1,\dots,w_{N_s}>0.
\]
Let \(\beta_s(\tau)\in\mathbb R^{2s}\) be the monomial basis vector from the previous phases, and define the row operator
\[
B^{(0)}(\tau):=\beta_s(\tau)^\top\in \mathbb R^{1\times 2s}.
\]
For segment \(i\) and sample index \(\nu\), define the sampled state
\begin{equation}
\label{eq:phase9_xinu}
x_{i,\nu}(c_i,T_i)
:=
B^{(0)}(\alpha_\nu T_i)c_i.
\end{equation}
Let
\[
\varphi_{\mathrm{obs}}:\mathbb R\to\mathbb R_{\ge 0}
\]
be any \(C^1\) smooth penalty function.
Then define
\begin{equation}
\label{eq:phase9_Kobs}
K_{\mathrm{soft\mbox{-}obs}}(c,T)
:=
\sum_{i=1}^{M}\sum_{\nu=1}^{N_s}
w_\nu \,\varphi_{\mathrm{obs}}(x_{i,\nu}(c_i,T_i)).
\end{equation}

\begin{remark}
\label{rem:phase9_obs_scope}
In the scalar canonical setting, \(\varphi_{\mathrm{obs}}\) should be interpreted as a minimal smooth spatial penalty.
The present phase does not attempt to model full geometric collision handling in \(\mathbb R^d\); that is deferred to Phase~10.
\end{remark}

\subsection{Differentiability of the minimal objective}
\label{subsec:phase9_diff}

The first task is to verify that the composite map
\[
\theta\mapsto c(\theta)\mapsto K(c(\theta),T)
\]
is differentiable.

\begin{proposition}[Differentiability of the raw Scheme~C objective]
\label{prop:phase9_differentiability}
Assume:
\begin{enumerate}
    \item the raw Scheme~C map \(c(\theta)=c^{\mathrm{C},2}(\theta)\) is \(C^1\) on \(\Theta_{\mathrm{adm}}(T_{\min})\);
    \item each segment Gram map \(T_i\mapsto H_i(T_i)\) is \(C^1\) on \((T_{\min},\infty)\);
    \item the obstacle penalty \(\varphi_{\mathrm{obs}}\) is \(C^1\).
\end{enumerate}
Then the objective \(J(\theta)\) defined by \eqref{eq:phase9_total_objective} is \(C^1\) on \(\Theta_{\mathrm{adm}}(T_{\min})\).
\end{proposition}

\begin{proof}
By assumption, \(c(\theta)\) is \(C^1\).
The control term \eqref{eq:phase9_Kctrl} is a finite sum of functions of the form
\[
(c_i,T_i)\mapsto \frac12 c_i^\top H_i(T_i)c_i,
\]
which are \(C^1\) because \(H_i\) is \(C^1\).
The time term \eqref{eq:phase9_Ktime} is linear in \(T\), hence \(C^\infty\).
For the obstacle term, each sampled state
\[
(c_i,T_i)\mapsto x_{i,\nu}(c_i,T_i)=B^{(0)}(\alpha_\nu T_i)c_i
\]
is \(C^1\), because \(B^{(0)}(\tau)\) is polynomial in \(\tau\).
Since \(\varphi_{\mathrm{obs}}\in C^1\), each summand
\[
(c_i,T_i)\mapsto w_\nu\varphi_{\mathrm{obs}}(x_{i,\nu}(c_i,T_i))
\]
is \(C^1\), and so is their finite sum.
Therefore \(K(c,T)\) is \(C^1\) in \((c,T)\), and by composition with the \(C^1\) map \(\theta\mapsto(c(\theta),T)\), the objective \(J(\theta)\) is \(C^1\).
\end{proof}

\subsection{Block derivatives of the objective with respect to \texorpdfstring{$c$}{c} and \texorpdfstring{$T$}{T}}
\label{subsec:phase9_block_derivatives}

We next compute the exact block gradients of \(K(c,T)\).

\subsubsection{Derivative with respect to the coefficient blocks}
\label{subsubsec:phase9_dKdc}

Define the block gradient
\[
g_c(c,T)
:=
\frac{\partial K}{\partial c}(c,T)
=
\begin{bmatrix}
g_{c,1}(c,T)\\
\vdots\\
g_{c,M}(c,T)
\end{bmatrix},
\qquad
g_{c,i}(c,T)\in\mathbb R^{2s}.
\]

\begin{proposition}[Exact formula for \texorpdfstring{$\partial K/\partial c$}{dK/dc}]
\label{prop:phase9_dKdc}
For each segment \(i\),
\begin{equation}
\label{eq:phase9_gc_block}
g_{c,i}(c,T)
=
H_i(T_i)c_i
+
\lambda_{\mathrm{obs}}
\sum_{\nu=1}^{N_s}
w_\nu\,
\varphi_{\mathrm{obs}}'(x_{i,\nu}(c_i,T_i))
\,B^{(0)}(\alpha_\nu T_i)^\top.
\end{equation}
\end{proposition}

\begin{proof}
Differentiate the control term:
\[
\frac{\partial}{\partial c_i}
\left(\frac12 c_i^\top H_i(T_i)c_i\right)
=
H_i(T_i)c_i,
\]
because \(H_i(T_i)\) is symmetric.
The time term does not depend on \(c_i\), so its derivative is zero.
For the obstacle term, each summand is
\[
w_\nu \varphi_{\mathrm{obs}}(x_{i,\nu}),
\qquad
x_{i,\nu}=B^{(0)}(\alpha_\nu T_i)c_i.
\]
By the scalar chain rule,
\[
\frac{\partial}{\partial c_i}
\bigl(w_\nu \varphi_{\mathrm{obs}}(x_{i,\nu})\bigr)
=
w_\nu \varphi_{\mathrm{obs}}'(x_{i,\nu})
\frac{\partial x_{i,\nu}}{\partial c_i}.
\]
Since \(x_{i,\nu}\) is linear in \(c_i\),
\[
\frac{\partial x_{i,\nu}}{\partial c_i}
=
B^{(0)}(\alpha_\nu T_i).
\]
As a column gradient, this contributes
\[
w_\nu \varphi_{\mathrm{obs}}'(x_{i,\nu})B^{(0)}(\alpha_\nu T_i)^\top.
\]
Summing over \(\nu\) and multiplying by \(\lambda_{\mathrm{obs}}\) gives \eqref{eq:phase9_gc_block}.
\end{proof}

\subsubsection{Derivative with respect to the time variables}
\label{subsubsec:phase9_dKdT}

Define
\[
g_T(c,T):=
\frac{\partial K}{\partial T}(c,T)
=
\begin{bmatrix}
g_{T,1}(c,T)\\
\vdots\\
g_{T,M}(c,T)
\end{bmatrix}
\in\mathbb R^M.
\]

Let
\[
B^{(1)}(\tau):=\frac{d}{d\tau}\beta_s(\tau)^\top\in\mathbb R^{1\times 2s}.
\]

\begin{proposition}[Exact formula for \texorpdfstring{$\partial K/\partial T$}{dK/dT}]
\label{prop:phase9_dKdT}
For each segment \(i\),
\begin{equation}
\label{eq:phase9_gT_block}
g_{T,i}(c,T)
=
\frac12 c_i^\top H_i'(T_i)c_i
+
\lambda_T
+
\lambda_{\mathrm{obs}}
\sum_{\nu=1}^{N_s}
w_\nu\,
\varphi_{\mathrm{obs}}'(x_{i,\nu}(c_i,T_i))
\,
\alpha_\nu B^{(1)}(\alpha_\nu T_i)c_i.
\end{equation}
\end{proposition}

\begin{proof}
Differentiate the control term:
\[
\frac{\partial}{\partial T_i}
\left(\frac12 c_i^\top H_i(T_i)c_i\right)
=
\frac12 c_i^\top H_i'(T_i)c_i.
\]
Differentiate the time term:
\[
\frac{\partial}{\partial T_i}\lambda_T\sum_{r=1}^{M}T_r=\lambda_T.
\]
For the obstacle term, each summand is
\[
w_\nu \varphi_{\mathrm{obs}}(x_{i,\nu}),
\qquad
x_{i,\nu}=B^{(0)}(\alpha_\nu T_i)c_i.
\]
Again by the chain rule,
\[
\frac{\partial}{\partial T_i}
\bigl(w_\nu \varphi_{\mathrm{obs}}(x_{i,\nu})\bigr)
=
w_\nu \varphi_{\mathrm{obs}}'(x_{i,\nu})
\frac{\partial x_{i,\nu}}{\partial T_i}.
\]
Because \(x_{i,\nu}=B^{(0)}(\tau)c_i\) with \(\tau=\alpha_\nu T_i\),
\[
\frac{\partial x_{i,\nu}}{\partial T_i}
=
\alpha_\nu B^{(1)}(\alpha_\nu T_i)c_i.
\]
Summing over \(\nu\) and adding all contributions yields \eqref{eq:phase9_gT_block}.
\end{proof}

\subsection{Exact reverse differentiation through the raw Scheme~C map}
\label{subsec:phase9_reverse_diff}

We now derive the exact chain rule for \(J(\theta)=K(c(\theta),T)\).

Write the Jacobians of the raw Scheme~C map as
\[
J_{c\bar q}(\theta)
:=
\frac{\partial c}{\partial \bar q}(\theta)
\in\mathbb R^{2sM\times(M-1)},
\qquad
J_{cT}(\theta)
:=
\frac{\partial c}{\partial T}(\theta)
\in\mathbb R^{2sM\times M}.
\]

\begin{theorem}[Exact chain rule for the raw Scheme~C optimization loop]
\label{thm:phase9_chain_rule}
Under the assumptions of Proposition~\ref{prop:phase9_differentiability},
the gradient of \(J(\theta)=K(c(\theta),T)\) is
\begin{align}
\label{eq:phase9_grad_q}
\nabla_{\bar q}J(\theta)
&=
J_{c\bar q}(\theta)^\top\, g_c(c(\theta),T),\\
\label{eq:phase9_grad_T}
\nabla_{T}J(\theta)
&=
J_{cT}(\theta)^\top\, g_c(c(\theta),T)
+
g_T(c(\theta),T).
\end{align}
Equivalently,
\begin{equation}
\label{eq:phase9_full_grad}
\nabla_\theta J(\theta)
=
\begin{bmatrix}
J_{c\bar q}^\top g_c\\
J_{cT}^\top g_c+g_T
\end{bmatrix}.
\end{equation}
\end{theorem}

\begin{proof}
Take an arbitrary differential \((d\bar q,dT)\).
By definition of the differential of a composite map,
\[
dJ
=
\left(\frac{\partial K}{\partial c}\right)^\top dc
+
\left(\frac{\partial K}{\partial T}\right)^\top dT.
\]
Because \(c=c(\bar q,T)\),
\[
dc
=
J_{c\bar q}\,d\bar q
+
J_{cT}\,dT.
\]
Substitute this into the previous expression:
\begin{align*}
dJ
&=
g_c^\top(J_{c\bar q}d\bar q + J_{cT}dT)
+
g_T^\top dT\\
&=
(J_{c\bar q}^\top g_c)^\top d\bar q
+
(J_{cT}^\top g_c + g_T)^\top dT.
\end{align*}
Since \(d\bar q\) and \(dT\) are arbitrary, the coefficients of these two differentials are exactly the gradients \(\nabla_{\bar q}J\) and \(\nabla_T J\), which proves \eqref{eq:phase9_grad_q}--\eqref{eq:phase9_grad_T}.
\end{proof}

\subsection{Exact block-banded backward accumulation}
\label{subsec:phase9_banded_accumulation}

The key computational advantage of raw Scheme~C is that the Jacobians \(J_{c\bar q}\) and \(J_{cT}\) are exactly banded, as proved in Phase~6.

In the canonical \(k=2\) setting, Phase~6 gives
\[
\frac{\partial c_i}{\partial q_j}=0
\qquad\text{for } j\notin\{i-2,i-1,i,i+1\},
\]
and
\[
\frac{\partial c_i}{\partial T_j}=0
\qquad\text{for } j\notin\{i-1,i,i+1\}.
\]

\begin{proposition}[Local backward accumulation formulas]
\label{prop:phase9_local_accumulation}
For each interior waypoint index \(j=1,\dots,M-1\),
\begin{equation}
\label{eq:phase9_local_accum_q}
\frac{\partial J}{\partial q_j}
=
\sum_{i=\max(1,j-1)}^{\min(M,j+2)}
\left(\frac{\partial c_i}{\partial q_j}\right)^\top g_{c,i}.
\end{equation}
For each time index \(j=1,\dots,M\),
\begin{equation}
\label{eq:phase9_local_accum_T}
\frac{\partial J}{\partial T_j}
=
g_{T,j}
+
\sum_{i=\max(1,j-1)}^{\min(M,j+1)}
\left(\frac{\partial c_i}{\partial T_j}\right)^\top g_{c,i}.
\end{equation}
\end{proposition}

\begin{proof}
From \eqref{eq:phase9_grad_q},
\[
\frac{\partial J}{\partial q_j}
=
\sum_{i=1}^{M}
\left(\frac{\partial c_i}{\partial q_j}\right)^\top g_{c,i}.
\]
By the exact local-support result of Phase~6,
\[
\frac{\partial c_i}{\partial q_j}=0
\quad\text{unless}\quad
j\in\{i-2,i-1,i,i+1\}.
\]
Equivalently,
\[
\frac{\partial c_i}{\partial q_j}=0
\quad\text{unless}\quad
i\in\{j-1,j,j+1,j+2\}.
\]
Therefore all other terms in the sum vanish, yielding \eqref{eq:phase9_local_accum_q}.

Likewise, from \eqref{eq:phase9_grad_T},
\[
\frac{\partial J}{\partial T_j}
=
g_{T,j}
+
\sum_{i=1}^{M}
\left(\frac{\partial c_i}{\partial T_j}\right)^\top g_{c,i}.
\]
Again by Phase~6,
\[
\frac{\partial c_i}{\partial T_j}=0
\quad\text{unless}\quad
j\in\{i-1,i,i+1\},
\]
equivalently
\[
i\in\{j-1,j,j+1\}.
\]
Hence all other terms vanish, and \eqref{eq:phase9_local_accum_T} follows.
\end{proof}

\begin{corollary}[Exact equivalence of dense and banded reverse accumulation]
\label{cor:phase9_dense_banded_equiv}
The gradient produced by the block-banded accumulation formulas
\eqref{eq:phase9_local_accum_q}--\eqref{eq:phase9_local_accum_T}
is exactly equal to the dense chain-rule gradient
\eqref{eq:phase9_full_grad}.
\end{corollary}

\begin{proof}
The banded formulas are obtained from the dense formulas by removing only those summands that are exactly zero by local support.
Therefore the two expressions are identical.
\end{proof}

\subsection{Complexity of the minimal differentiable loop}
\label{subsec:phase9_complexity}

We now state the exact complexity result for the \emph{canonical minimal loop} of this phase.
Because \(s=4\), \(k=2\), and the number of obstacle samples \(N_s\) is fixed, all per-segment operations are \(O(1)\).

\begin{theorem}[Linear-time complexity of the canonical minimal loop]
\label{thm:phase9_linear_complexity}
Assume:
\begin{enumerate}
    \item \(s=4\), \(k=2\), and \(N_s\) is fixed;
    \item each local coefficient block \(c_i(\theta)\) and its nonzero Jacobian blocks
    \[
    \frac{\partial c_i}{\partial q_j},\qquad
    \frac{\partial c_i}{\partial T_j}
    \]
    are available in \(O(1)\) time per admissible index pair;
    \item the obstacle penalty \(\varphi_{\mathrm{obs}}\) and its derivative are evaluable in \(O(1)\) time.
\end{enumerate}
Then one evaluation of:
\begin{enumerate}
    \item the objective \(J(\theta)\),
    \item the block gradients \(g_c(c(\theta),T)\) and \(g_T(c(\theta),T)\),
    \item and the full parameter gradient \(\nabla_\theta J(\theta)\)
\end{enumerate}
can be computed in
\begin{equation}
\label{eq:phase9_linear_complexity}
O(M)
\end{equation}
time and \(O(M)\) memory.
\end{theorem}

\begin{proof}
We count the operations component by component.

\paragraph{Step 1: coefficient evaluation.}
Since the raw Scheme~C map is local and \(k=2\) is fixed, each segment coefficient \(c_i(\theta)\in\mathbb R^8\) is computed from a fixed-size local system, hence in \(O(1)\) time.
Over all \(M\) segments, this costs \(O(M)\).

\paragraph{Step 2: objective evaluation.}
For each segment, the control term \(\frac12 c_i^\top H_i(T_i)c_i\) requires \(O(1)\) operations because the block size is fixed.
The time term contributes \(O(1)\) per segment.
The obstacle term uses \(N_s\) fixed samples per segment, hence also \(O(1)\) per segment.
Therefore the total objective evaluation costs \(O(M)\).

\paragraph{Step 3: block gradients with respect to \(c\) and \(T\).}
By \eqref{eq:phase9_gc_block} and \eqref{eq:phase9_gT_block}, each block \(g_{c,i}\) and scalar \(g_{T,i}\) is computed from a fixed number of fixed-size matrix-vector and scalar operations.
Thus this step is also \(O(M)\).

\paragraph{Step 4: reverse accumulation to \(\bar q\) and \(T\).}
For each waypoint index \(j\), formula \eqref{eq:phase9_local_accum_q} sums at most four nonzero contributions.
For each time index \(j\), formula \eqref{eq:phase9_local_accum_T} sums at most three nonzero contributions plus one local term \(g_{T,j}\).
Hence each gradient component is accumulated in \(O(1)\), and the total reverse sweep is \(O(M)\).

\paragraph{Step 5: memory.}
All block quantities
\[
c_i,\quad g_{c,i},\quad g_{T,i},
\]
and the gradient vectors have size \(O(M)\), while each block has fixed dimension.
Thus the memory footprint is \(O(M)\).

Combining all steps proves \eqref{eq:phase9_linear_complexity}.
\end{proof}

\begin{remark}[Marked limitation]
\label{rem:phase9_complexity_limitation}
The full \(O(k^2M)\) complexity theorem for general \(k\), together with a complete block-banded implementation theorem for the full BLOM framework, is \textbf{not completed in this phase}.
That statement belongs to Phase~10.
Phase~9 proves only the canonical linear-time result for the minimal raw Scheme~C loop.
\end{remark}

\subsection{A rigorous local descent guarantee}
\label{subsec:phase9_descent}

The goal of Phase~9 is not to prove global convergence of a nonconvex planner.
It is sufficient to prove that, once the exact gradient is available, the minimal loop admits a valid descent step.

\begin{theorem}[Local descent under exact gradient and smoothness]
\label{thm:phase9_descent}
Assume that \(J\in C^1(\Theta_{\mathrm{adm}}(T_{\min}))\) and that \(\nabla J\) is Lipschitz continuous with constant \(L>0\) on a convex neighborhood \(U\subset\Theta_{\mathrm{adm}}(T_{\min})\) of a point \(\theta\in U\).
Define the negative-gradient direction
\[
d(\theta):=-\nabla J(\theta).
\]
Let
\begin{equation}
\label{eq:phase9_alpha_feas}
\alpha_{\mathrm{feas}}(\theta)
:=
\inf\left\{
\alpha>0:
T_i-\alpha \nabla_{T_i}J(\theta)\le T_{\min}
\text{ for some } i
\right\},
\end{equation}
with the convention \(\alpha_{\mathrm{feas}}(\theta)=+\infty\) if the set is empty.
Then for every step size
\[
0<\alpha<\min\{\alpha_{\mathrm{feas}}(\theta),\,1/L\},
\]
the point
\[
\theta^+ := \theta-\alpha \nabla J(\theta)
\]
remains in \(\Theta_{\mathrm{adm}}(T_{\min})\) and satisfies
\begin{equation}
\label{eq:phase9_descent_ineq}
J(\theta^+)
\le
J(\theta)
-
\frac{\alpha}{2}\|\nabla J(\theta)\|_2^2.
\end{equation}
\end{theorem}

\begin{proof}
By the definition of \(\alpha_{\mathrm{feas}}(\theta)\), any \(0<\alpha<\alpha_{\mathrm{feas}}(\theta)\) guarantees
\[
T_i-\alpha \nabla_{T_i}J(\theta)>T_{\min}
\qquad
\forall i,
\]
so \(\theta^+\in \Theta_{\mathrm{adm}}(T_{\min})\).

Because \(\nabla J\) is Lipschitz with constant \(L\) on \(U\), the descent lemma gives
\[
J(\theta^+)
\le
J(\theta)
+
\nabla J(\theta)^\top(\theta^+-\theta)
+
\frac{L}{2}\|\theta^+-\theta\|_2^2.
\]
Since \(\theta^+-\theta=-\alpha \nabla J(\theta)\),
\[
J(\theta^+)
\le
J(\theta)
-
\alpha \|\nabla J(\theta)\|_2^2
+
\frac{L}{2}\alpha^2 \|\nabla J(\theta)\|_2^2.
\]
Factor out the gradient norm:
\[
J(\theta^+)
\le
J(\theta)
-
\alpha\left(1-\frac{L\alpha}{2}\right)\|\nabla J(\theta)\|_2^2.
\]
If \(\alpha<1/L\), then
\[
1-\frac{L\alpha}{2}>\frac12,
\]
hence
\[
J(\theta^+)
\le
J(\theta)
-
\frac{\alpha}{2}\|\nabla J(\theta)\|_2^2.
\]
This proves \eqref{eq:phase9_descent_ineq}.
\end{proof}

\begin{remark}[Marked limitation]
\label{rem:phase9_global_opt_limitation}
A global convergence theorem for repeated optimization iterates in the full nonconvex setting is \textbf{not completed at the present stage}.
Phase~9 proves only that the exact differentiable loop admits valid descent steps, which is the mathematically correct level of guarantee for a minimal optimization phase.
\end{remark}

\subsection{Dense checker versus banded implementation}
\label{subsec:phase9_dense_checker}

For practical verification, it is useful to distinguish between:
\begin{enumerate}
    \item a dense checker implementation, which forms the full Jacobians \(J_{c\bar q}\) and \(J_{cT}\) explicitly;
    \item the true block-banded implementation, which only visits the nonzero Jacobian blocks from Proposition~\ref{prop:phase9_local_accumulation}.
\end{enumerate}

\begin{proposition}[Dense checker and banded implementation are mathematically identical]
\label{prop:phase9_dense_checker}
Suppose the dense checker uses the exact Jacobians \(J_{c\bar q}\) and \(J_{cT}\), while the banded implementation uses only the nonzero block entries guaranteed by raw Scheme~C local support.
Then both implementations produce exactly the same value of \(\nabla_\theta J(\theta)\).
\end{proposition}

\begin{proof}
This is an immediate consequence of Corollary~\ref{cor:phase9_dense_banded_equiv}.
The banded implementation removes only entries that are exactly zero.
Therefore the accumulated sums are identical to the dense chain-rule expressions.
\end{proof}

\subsection{What is proved and what remains open in Phase~9}
\label{subsec:phase9_status}

We summarize the exact status of this phase.

\paragraph{Fully proved in this chapter.}
\begin{enumerate}
    \item the differentiability of the minimal raw Scheme~C objective;
    \item the exact formulas for \(\partial K/\partial c\) and \(\partial K/\partial T\);
    \item the exact reverse-differentiation chain rule
    \[
    \nabla_\theta J
    =
    \begin{bmatrix}
    J_{c\bar q}^\top g_c\\
    J_{cT}^\top g_c+g_T
    \end{bmatrix};
    \]
    \item the local block-banded accumulation formulas for \(\nabla_{\bar q}J\) and \(\nabla_TJ\);
    \item the exact equivalence of dense and banded reverse accumulation;
    \item the \(O(M)\) complexity theorem for the canonical minimal loop;
    \item the existence of a valid local descent step under standard smoothness assumptions.
\end{enumerate}

\paragraph{Not completed at the present stage.}
\begin{enumerate}
    \item the full \(O(k^2M)\) complexity theorem for general \(k\);
    \item the complete block-banded theorem for the full BLOM family;
    \item the integration of hard dynamics constraints, ESDF-heavy collision systems, and full planner machinery;
    \item a global convergence theorem for the nonconvex optimization process.
\end{enumerate}

Accordingly, the mathematically correct summary of Phase~9 is:
\[
\boxed{
\text{raw Scheme~C is already a valid differentiable representation with an exact banded backward pass,}
}
\]
and
\[
\boxed{
\text{the minimal optimization loop is rigorous enough to justify Phase~10, but not yet a full planning theory.}
}